\begin{table}[H]
\centering
\caption{Cohorts contributing to each event time in the headline event study}
\label{tab:cohort_contributions}
% latex table generated in R 4.6.0 by xtable 1.8-8 package
% Thu Sep 17 13:21:23 2026
\begingroup\small
\adjustbox{max width=\textwidth}{%
\begin{tabular}{llll}
  \toprule
Event time $e$ & Contributing cohorts & Cohorts & Treated recipients \\ 
  \midrule
-5 & 2021, 2022, 2023, 2024 & 4 & 40 \\ 
  -4 & 2021, 2022, 2023, 2024 & 4 & 40 \\ 
  -3 & 2021, 2022, 2023, 2024 & 4 & 40 \\ 
  -2 & 2021, 2022, 2023, 2024 & 4 & 40 \\ 
  -1 & --- & 0 & 0 \\ 
  0 & 2021, 2022, 2023, 2024 & 4 & 40 \\ 
  1 & 2021, 2022, 2023 & 3 & 32 \\ 
  2 & 2021, 2022 & 2 & 21 \\ 
  3 & 2021 & 1 & 11 \\ 
   \bottomrule
\end{tabular}
}
\endgroup
\par\vspace{4pt}
\begin{minipage}{\linewidth}\footnotesize
Notes: Composition of the dynamic aggregation used by Figure~\ref{fig:did_combined_es} and the HonestDiD analysis. A cohort contributes to event time $e$ when its $ATT(g, g+e)$ cell is estimable in the stored headline fit (non-missing, $se > 0$); $e = -1$ is the base period and contributes none. Counts are treated recipients of the contributing cohorts.\par
Source: OECD CRS (Rio adaptation markers); UNFCCC NAP Central.
\end{minipage}
\end{table}
